Tato sekce nabízí krátkou sbírku praktických šablon a ukázek pro redesign předmětů s AI‑aware komponentami; upravte podle kurzu a pravidel instituce.

\bigskip

\textbf{Krátký text do sylabu (vzor)}
\begin{quote}
Použití generativní AI: AI nástroje povoleny pouze pokud student při odevzdání přiloží „AI log“ a stručnou reflexi. Uveďte: (1) nástroj a verzi, (2) přesné prompty/vstupy, (3) co AI vygenerovala a jak jste výstup upravili, (4) krátkou reflexi o přínosu a limitech. Neúplná deklarace může snížit hodnocení.
\end{quote}

\bigskip

\textbf{Milníky (template)}
\begin{enumerate}
  \item M1 Návrh (týd.2): cíl, postup, zdroje. Hodnocení: jasnost plánu.
  \item M2 Prototyp (týd.4): funkční část, logy včetně AI logu. Hodnocení: reprodukovatelnost.
  \item M3 Testy (týd.8): metriky, ověření dat. M4 Závěr (týd.12): finální artefakt + 1–2 s. reflexe.
\end{enumerate}

\bigskip

\textbf{Prompt‑log (šablona)}\\
\begin{PromptBlock}

- Nástroj/verze; - Přesný prompt; - Co AI vygenerovala; - Jak jsem upravil/ověřil; - Reflexe (2–4 odstavce).
\end{PromptBlock}


\bigskip

\textbf{Rubrika a checklist pro lektora}\\
Rubrika: Proces 40\% (kroky, zdůvodnění, logy), Reprodukovatelnost 25\% (data, metadata, hashe), Reflexe AI 20\%, Konečný artefakt 15\%.\\
Checklist: vymezit proces vs produkt; aktualizovat sylabus; navrhnout 3–4 milníky; zajistit technické požadavky (metadata, formát logu, volitelné hashování); provést pilot.